\chapter{Complex Scalar Field and Internal Symmetry}

\begin{remarkbox}
A complex scalar field is the simplest field theory with a nontrivial internal
symmetry. It can be understood as two real scalar fields, but its complex form
makes a global \(U(1)\) phase symmetry manifest. Noether's theorem then gives a
conserved current and charge. By promoting the phase symmetry from global to
local, one is naturally led to the gauge covariant derivative and to coupling
with electromagnetism.
\end{remarkbox}

\section{Complex Field as Two Real Fields}

A complex scalar field is a function
\[
    \phi : M \to \mathbb{C}.
\]
It may be decomposed into two real fields,
\[
    \phi=\frac{1}{\sqrt{2}}(\phi_1+i\phi_2),
    \qquad
    \phi^*=\frac{1}{\sqrt{2}}(\phi_1-i\phi_2),
\]
where \(\phi_1\) and \(\phi_2\) are real scalar fields. The factor
\(1/\sqrt{2}\) is conventional and gives convenient normalizations.

The free complex scalar field is described by the Lagrangian density
\[
    \mathcal{L}
    =\partial_\mu\phi^*\partial^\mu\phi-m^2\phi^*\phi.
\]
In terms of \(\phi_1\) and \(\phi_2\), this becomes
\[
    \mathcal{L}
    =\frac{1}{2}\partial_\mu\phi_1\partial^\mu\phi_1
     +\frac{1}{2}\partial_\mu\phi_2\partial^\mu\phi_2
     -\frac{1}{2}m^2\phi_1^2
     -\frac{1}{2}m^2\phi_2^2.
\]
Thus a free complex scalar field is equivalent to two real scalar fields with
the same mass.

\begin{definitionbox}
A \emph{complex scalar field} is a scalar field with complex values. As a real
field theory it contains two real scalar degrees of freedom, which may be taken
as the real and imaginary parts of the field.
\end{definitionbox}

The complex notation is not merely aesthetic. It makes the internal rotation
between \(\phi_1\) and \(\phi_2\) transparent. This internal rotation is the
origin of the conserved \(U(1)\) charge.

The Euler--Lagrange equations may be obtained by treating \(\phi\) and \(\phi^*\)
as independent variables during the variation. Varying with respect to
\(\phi^*\) gives
\[
    (\Box+m^2)\phi=0.
\]
Varying with respect to \(\phi\) gives the complex conjugate equation
\[
    (\Box+m^2)\phi^*=0.
\]
Equivalently, \(\phi_1\) and \(\phi_2\) each satisfy the real Klein--Gordon
equation.

\section{Global \texorpdfstring{\(U(1)\)}{U(1)} Symmetry}

The free complex scalar Lagrangian is invariant under constant phase rotations
\[
    \phi\longrightarrow e^{i\alpha}\phi,
    \qquad
    \phi^*\longrightarrow e^{-i\alpha}\phi^*,
\]
where \(\alpha\) is a real constant. This is a global \(U(1)\) symmetry.

For infinitesimal \(\alpha\),
\[
    e^{i\alpha}=1+i\alpha+O(\alpha^2),
\]
so the field variations are
\[
    \delta\phi=i\alpha\phi,
    \qquad
    \delta\phi^*=-i\alpha\phi^*.
\]
The derivatives transform in the same way:
\[
    \delta(\partial_\mu\phi)=i\alpha\partial_\mu\phi,
    \qquad
    \delta(\partial_\mu\phi^*)=-i\alpha\partial_\mu\phi^*.
\]
Substituting these variations into
\[
    \mathcal{L}
    =\partial_\mu\phi^*\partial^\mu\phi-m^2\phi^*\phi
\]
shows that
\[
    \delta\mathcal{L}=0.
\]
Thus the symmetry is exact and leaves the Lagrangian density itself invariant,
not merely up to a total derivative.

\begin{definitionbox}
The group \(U(1)\) is the group of complex phases \(e^{i\alpha}\). A global
\(U(1)\) symmetry acts by multiplying the field by a constant phase at every
spacetime point.
\end{definitionbox}

In terms of the two real fields, this symmetry is an internal rotation:
\[
    \begin{pmatrix}\phi_1\\ \phi_2\end{pmatrix}
    \longrightarrow
    \begin{pmatrix}
    \cos\alpha & -\sin\alpha\\
    \sin\alpha & \cos\alpha
    \end{pmatrix}
    \begin{pmatrix}\phi_1\\ \phi_2\end{pmatrix}.
\]
It is called internal because it rotates field components into one another
without rotating physical space.

\section{Conserved Current and Charge}

Noether's theorem associates a conserved current with the global \(U(1)\)
symmetry. Since
\[
    \mathcal{L}=\partial_\mu\phi^*\partial^\mu\phi-m^2\phi^*\phi,
\]
we have
\[
    \frac{\partial\mathcal{L}}{\partial(\partial_\mu\phi)}=\partial^\mu\phi^*,
    \qquad
    \frac{\partial\mathcal{L}}{\partial(\partial_\mu\phi^*)}=\partial^\mu\phi.
\]
Using the infinitesimal variations without the common parameter \(\alpha\),
\[
    \Delta\phi=i\phi,
    \qquad
    \Delta\phi^*=-i\phi^*,
\]
the Noether current is
\[
    J^\mu
    =\frac{\partial\mathcal{L}}{\partial(\partial_\mu\phi)}\Delta\phi
     +\frac{\partial\mathcal{L}}{\partial(\partial_\mu\phi^*)}\Delta\phi^*.
\]
Therefore
\[
    J^\mu=i\phi\partial^\mu\phi^*-i\phi^*\partial^\mu\phi.
\]
Equivalently,
\[
    J^\mu=i(\phi\partial^\mu\phi^*-\phi^*\partial^\mu\phi).
\]
Some authors use the opposite sign convention. The convention is not important
as long as the definition of charge is used consistently.

On shell,
\[
    \partial_\mu J^\mu=0.
\]
The conserved charge is
\[
    Q=\int d^d x\,J^0.
\]
Using \(\partial^0=\partial_t\), one has
\[
    J^0=i(\phi\dot\phi^*-\phi^*\dot\phi).
\]
Thus
\[
    Q=i\int d^d x\,(\phi\dot\phi^*-\phi^*\dot\phi).
\]

\begin{theorembox}
The global phase symmetry of the complex scalar field gives the conserved
Noether current
\[
    J^\mu=i(\phi\partial^\mu\phi^*-\phi^*\partial^\mu\phi),
    \qquad
    \partial_\mu J^\mu=0
\]
on solutions of the field equations.
\end{theorembox}

This current is the classical precursor of charge current for a charged scalar
field. In a quantum theory, the corresponding charge counts particles and
antiparticles with opposite signs. In the classical theory, it measures the
oriented internal rotation carried by the field configuration.

\section{Coupling to an External Electromagnetic Potential}

A charged scalar field can be coupled to a prescribed electromagnetic potential
\(A_\mu(x)\). At first, we treat \(A_\mu\) as an external background, not as a
dynamical field. The guiding principle is that ordinary derivatives should be
replaced by covariant derivatives:
\[
    \partial_\mu\phi\longrightarrow D_\mu\phi.
\]
For a field of charge \(q\), a common convention is
\[
    D_\mu\phi=(\partial_\mu+iqA_\mu)\phi.
\]
The complex conjugate field transforms with the opposite charge, so
\[
    D_\mu\phi^*=(\partial_\mu-iqA_\mu)\phi^*.
\]

The Lagrangian density becomes
\[
    \mathcal{L}
    =(D_\mu\phi)^*D^\mu\phi-m^2\phi^*\phi.
\]
Here
\[
    (D_\mu\phi)^*=D_\mu\phi^*
\]
with the above convention. Expanding the covariant derivative shows that the
background field couples to the scalar current and also introduces a term
quadratic in \(A_\mu\).

\begin{remarkbox}
The electromagnetic potential does not couple to a real scalar field in the same
minimal way, because a real scalar field has no continuous phase that can carry
\(U(1)\) charge. Charged scalar matter is naturally complex.
\end{remarkbox}

If \(A_\mu\) is fixed externally, the scalar field evolves in a prescribed
electromagnetic background. If \(A_\mu\) is dynamical, one must add the Maxwell
term to the action and vary with respect to both \(\phi\) and \(A_\mu\).

\section{Local Phase Symmetry}

The global phase symmetry uses a constant parameter \(\alpha\). A local phase
transformation allows the parameter to depend on spacetime:
\[
    \phi(x)\longrightarrow \phi'(x)=e^{i\alpha(x)}\phi(x).
\]
If one uses ordinary derivatives, the derivative transforms as
\[
    \partial_\mu\phi'(x)
    =e^{i\alpha(x)}\partial_\mu\phi(x)
     +i(\partial_\mu\alpha)e^{i\alpha(x)}\phi(x).
\]
The second term prevents \(\partial_\mu\phi\) from transforming in the same way
as \(\phi\). Therefore the free Lagrangian with ordinary derivatives is not
invariant under local phase transformations.

The remedy is to introduce a gauge field \(A_\mu\) whose transformation cancels
the extra derivative of \(\alpha\). With the convention
\[
    D_\mu=\partial_\mu+iqA_\mu,
\]
one takes
\[
    A_\mu\longrightarrow A'_\mu=A_\mu-\frac{1}{q}\partial_\mu\alpha.
\]
Then
\[
    D'_\mu\phi'=e^{i\alpha}D_\mu\phi.
\]
This is the defining covariance property of the gauge covariant derivative.

\begin{definitionbox}
A \emph{local phase symmetry} is a transformation with spacetime-dependent phase
\(\alpha(x)\). It requires a gauge field because ordinary derivatives of the
field do not transform covariantly under local phase rotations.
\end{definitionbox}

This is the conceptual origin of gauge theory: demanding local internal symmetry
forces the introduction of a compensating connection-like field.

\section{Minimal Coupling}

Minimal coupling is the prescription
\[
    \partial_\mu\longrightarrow D_\mu=\partial_\mu+iqA_\mu.
\]
Applied to the free complex scalar field, it gives
\[
    \mathcal{L}_{\mathrm{matter}}
    =(D_\mu\phi)^*D^\mu\phi-m^2\phi^*\phi.
\]
If the electromagnetic field is dynamical, the full scalar electrodynamics
Lagrangian is
\[
    \mathcal{L}
    =-\frac{1}{4}F_{\mu\nu}F^{\mu\nu}
     +(D_\mu\phi)^*D^\mu\phi-m^2\phi^*\phi,
\]
where
\[
    F_{\mu\nu}=\partial_\mu A_\nu-\partial_\nu A_\mu.
\]
This theory is invariant under
\[
    \phi\longrightarrow e^{i\alpha(x)}\phi,
    \qquad
    A_\mu\longrightarrow A_\mu-\frac{1}{q}\partial_\mu\alpha.
\]

\begin{examplebox}
The replacement \(\partial_\mu\to D_\mu\) is not an arbitrary trick. It is forced
by the requirement that derivatives of charged fields transform covariantly
under local phase rotations.
\end{examplebox}

The word minimal means that one introduces the gauge field through the derivative
and does not add extra higher-order or nonminimal terms such as
\(F_{\mu\nu}F^{\mu\nu}\phi^*\phi\). Such terms can be considered in effective
field theories, but they are not part of the minimal coupling prescription.

\section{Gauge Covariant Derivative}

The gauge covariant derivative is designed to transform like the field itself:
\[
    D'_\mu\phi'=e^{i\alpha}D_\mu\phi.
\]
This property ensures that
\[
    (D_\mu\phi)^*D^\mu\phi
\]
is invariant under local \(U(1)\) transformations. The mass term
\[
    m^2\phi^*\phi
\]
is also invariant, because the phase cancels between \(\phi\) and \(\phi^*\).

The commutator of two covariant derivatives gives the field strength. With
\[
    D_\mu=\partial_\mu+iqA_\mu,
\]
one finds
\[
    [D_\mu,D_\nu]\phi=iqF_{\mu\nu}\phi.
\]
Thus the electromagnetic field strength measures the noncommutativity of gauge
covariant derivatives.

\begin{definitionbox}
For a charged scalar field, the gauge covariant derivative is
\[
    D_\mu\phi=(\partial_\mu+iqA_\mu)\phi.
\]
Its commutator satisfies
\[
    [D_\mu,D_\nu]\phi=iqF_{\mu\nu}\phi.
\]
\end{definitionbox}

This formula is the Abelian prototype of Yang--Mills theory. In non-Abelian
gauge theory, the gauge potential and the field strength take values in a Lie
algebra, and the commutator term becomes structurally essential.

\section{Physical Meaning of Charge}

In the complex scalar theory, charge is the Noether charge associated with
internal phase rotations. The field has a two-dimensional internal plane
spanned by \(\phi_1\) and \(\phi_2\). A charged configuration carries a sense of
rotation in that internal plane.

Using
\[
    \phi=\frac{1}{\sqrt{2}}(\phi_1+i\phi_2),
\]
the charge density can be written in terms of real fields as
\[
    J^0=\phi_2\dot\phi_1-\phi_1\dot\phi_2,
\]
up to the overall sign convention chosen for \(J^\mu\). This is analogous to
angular momentum in the internal \((\phi_1,\phi_2)\)-plane.

\begin{remarkbox}
A complex scalar field is not ``more mysterious'' than two real scalar fields.
What is new is the internal rotational symmetry between them. The conserved
charge is the Noether charge of that internal rotation.
\end{remarkbox}

When the field is coupled to electromagnetism, this internal charge becomes the
source of the electromagnetic field. In scalar electrodynamics, varying the
action with respect to \(A_\mu\) gives Maxwell equations with a matter current
constructed from \(\phi\) and \(D_\mu\phi\).

\section{Charged Scalar Field as a Classical Field Theory}

The charged scalar field provides the simplest classical model of matter coupled
to a gauge field. Its dynamical variables are
\[
    \phi,\qquad \phi^*,\qquad A_\mu.
\]
The action is
\[
    S=\int d^{d+1}x\left[
        -\frac{1}{4}F_{\mu\nu}F^{\mu\nu}
        +(D_\mu\phi)^*D^\mu\phi
        -m^2\phi^*\phi
    \right].
\]
The scalar field equation is
\[
    D_\mu D^\mu\phi+m^2\phi=0,
\]
with the complex conjugate equation for \(\phi^*\). The gauge field equation has
the form
\[
    \partial_\mu F^{\mu\nu}=j^\nu,
\]
where the matter current is obtained by varying the matter action with respect
to \(A_\nu\).

This theory illustrates several themes that will recur in the next chapters:
\begin{enumerate}
    \item internal symmetry produces charge;
    \item local symmetry requires a gauge field;
    \item the gauge potential is not itself gauge invariant;
    \item the field strength is gauge invariant in the Abelian theory;
    \item the covariant derivative is the correct derivative for charged fields.
\end{enumerate}

\begin{summarybox}
The complex scalar field is equivalent to two real scalar fields but displays a
natural global \(U(1)\) symmetry. Noether's theorem gives a conserved current
and charge. Promoting the phase symmetry to a local symmetry forces the
introduction of a gauge field and the covariant derivative
\(D_\mu=\partial_\mu+iqA_\mu\). This is the simplest route from scalar field
theory to electromagnetism and, ultimately, to Yang--Mills gauge theory.
\end{summarybox}
